\section{Algebraic and kinematic closure on the solid frame}
\label{sec:algebraic_kinematic_closure}

The weak forms in Section~\ref{sec:reference_solid_weak_forms} are not a
closed finite-element system unless the partition constraints, thermodynamic
multiplier relations, transfer-potential evolution laws, and phase kinematics
are solved or supplied as material equations.  This section records those
relations in the variables needed by a solid-reference implementation.

\subsection{Partition and storage equations}
\label{sec:partition_storage_equations}

For each phase \(\xi\) and component \(\alpha\), the component partial density,
phase partial density, intrinsic density, volume fraction, and mass fraction
are tied by
\begin{subequations}
	\label{eq:implementation_partition_relations}
	\begin{align}
		\rho_\xi^\alpha
		 & =\phi_\xi\bar\rho_\xi\eta_\xi^\alpha,
		\label{eq:implementation_component_density_partition} \\
		\rho_\xi
		 & =\phi_\xi\bar\rho_\xi,
		\label{eq:implementation_phase_density_partition} \\
		\sum_\xi\phi_\xi
		 & =1,
		\label{eq:implementation_volume_partition} \\
		\sum_\alpha\eta_\xi^\alpha
		 & =1 .
		\label{eq:implementation_mass_fraction_partition}
	\end{align}
\end{subequations}
The phase-attached material mass identity gives the pointwise storage relation
\begin{equation}
	\rho_{\xi0}^\alpha+\mc C_\xi^\alpha
	=
	J_\xi\phi_\xi\bar\rho_\xi\eta_\xi^\alpha .
	\label{eq:implementation_phase_material_storage}
\end{equation}
For a fluid phase \(f\), the solid-reference storage that enters the component
balance is \(J\phi_f\bar\rho_f\eta_f^\alpha\), where \(J\) is the solid
skeleton Jacobian.  Equation~\eqref{eq:implementation_phase_material_storage}
is still needed because the current source rate is
\begin{equation}
	\dot c_\xi^\alpha
	=
	\frac{\dot{\mc C}_\xi^\alpha}{J_\xi}
	=
	\sum_m\nu_{\xi (m)}^\alpha\dot r_{(m)} .
	\label{eq:implementation_conversion_rate_constraint}
\end{equation}
Thus \(J_\xi\) is part of the phase-attached definition of
\(\dot{\mc C}_\xi^\alpha\).  The solid-reference implementation keeps this
distinction explicit.  The input deck registers a finite phase set \(\mc P\)
and selects its reference solid phase \(s\in\mc P\).  The weak forms are pulled
back with that solid skeleton map, while every registered mobile phase
\(a\in\mc P\setminus\{s\}\) may carry phase-history deformation gradients and
Jacobians when a closure needs phase-attached material rates.

\subsection{Fluid material derivatives on the solid frame}
\label{sec:fluid_material_derivatives_solid_frame}

For each registered mobile phase \(a\in\mc P\setminus\{s\}\), the
solid-frame convective velocity is reconstructed from the same reference
relative mass flux that enters the component balance:
\begin{equation}
	\mbf v_a
	=
	\mbf v_s+\frac{1}{J_s\rho_a}\mbf F_s\mbf W_a,
	\qquad
	\mbf c_a
	=
	\mbf F_s^{-1}\left(\mbf v_a-\mbf v_s\right)
	=
	\frac{1}{J_s\rho_a}\mbf W_a .
	\label{eq:implementation_phase_velocity_reconstruction}
\end{equation}
Here \(\rho_a=\phi_a\bar\rho_a\) is the current bulk phase density per unit
current mixture volume; using the intrinsic density \(\bar\rho_a\) alone would
omit the phase volume fraction and give an incorrect material derivative.
Therefore any phase-attached scalar or tensor field \(q_a\) stored on the solid
mesh has
\begin{equation}
	\frac{D_a q_a}{Dt}
	=
	\left.\frac{\partial q_a}{\partial t}\right|_{\mbf X}
	+\mbf c_a\cdot\Grad q_a .
	\label{eq:fluid_material_derivative_on_solid_frame}
\end{equation}
The phase-history kinematics use each mobile phase velocity gradient in the
current configuration,
\begin{equation}
	\mbf l_a
	=
	\grad \mbf v_a
	=
	\Grad\mbf v_a\,\mbf F_s^{-1},
	\label{eq:implementation_phase_velocity_gradient}
\end{equation}
and evolve on the solid mesh according to
\begin{subequations}
	\label{eq:implementation_phase_history_kinematics}
	\begin{align}
		\left.\frac{\partial \mbf F_a}{\partial t}\right|_{\mbf X}
		+\mbf c_a\cdot\Grad\mbf F_a
		 & =
		\mbf l_a\mbf F_a,
		\label{eq:implementation_phase_deformation_gradient_history} \\
		\left.\frac{\partial J_a}{\partial t}\right|_{\mbf X}
		+\mbf c_a\cdot\Grad J_a
		 & =
		J_a\tr\mbf l_a .
		\label{eq:implementation_phase_jacobian_history}
	\end{align}
\end{subequations}
Here \(\mbf F_a\) and \(J_a\) are phase-history kinematics associated with the
phase labels, stored on the solid mesh.  They are distinct from
\(\mbf F_s\) and \(J_s\), which define the reference configuration, all
Grad/Div operators, boundary measures, and Piola transformations in the weak
forms.  Input decks must initialize a history-free phase with
\(\mbf F_a=\mbf I\) and \(J_a=1\), or provide another mutually consistent
history.  If a phase is inactive, the implementation sets both
\(\mbf c_a=\mbf 0\) and \(\mbf l_a=\mbf 0\), thereby preserving the stored
history until a documented closure reactivates that phase with consistent data.

\subsection{Thermodynamic multiplier and partition relations}
\label{sec:thermodynamic_partition_relations}

The general EOS input is a user-supplied Helmholtz free-energy density per unit
current phase volume,
\(\mc A_a=\mc A_a(\rho_a^1,\ldots,\rho_a^N,\theta,\ldots)\), where the
\(\rho_a^\alpha\) are constituent mass densities per unit current phase
volume.  The parsed material differentiates this single potential rather than
accepting pressure or chemical potentials as independent constitutive inputs:
\begin{subequations}
\begin{align}
 \mu_a^\alpha &=
 \left.\frac{\partial\mc A_a}{\partial\rho_a^\alpha}\right|_{\rho_a^{\beta\ne\alpha},\theta,\ldots},
 \label{eq:implementation_eos_chemical_potential}\\
 \bar p_a &= \sum_\alpha\rho_a^\alpha\mu_a^\alpha-\mc A_a,
 \label{eq:implementation_eos_pressure}\\
 s_a &= -\left.\frac{\partial\mc A_a}{\partial\theta}\right|_{\rho_a^1,\ldots,\rho_a^N,\ldots}.
 \label{eq:implementation_eos_entropy}
\end{align}
\end{subequations}
Second derivatives of \(\mc A_a\) provide the composition and thermal
thermodynamic Jacobian.  Thus pressure, chemical potentials, entropy, and their
Newton couplings remain mutually consistent for every admissible user
expression.  For a single density argument, this Euler relation reduces to
\begin{equation}
	\bar p_\xi
	=
	\bar\rho_\xi^2
	\left.\frac{\partial\psi_\xi}{\partial\bar\rho_\xi}\right|_{\mbf F_\xi,\eta_\xi^\beta,\theta}.
	\label{eq:implementation_phase_pressure}
\end{equation}
The equivalent pore pressure is
\begin{equation}
	\bar p_E
	=
	-\lambda
	=
	\sum_{f\in\mc F}S_f\bar p_f-\gamma .
	\label{eq:implementation_equivalent_pressure_relation}
\end{equation}
The individual capillary relation
\(\lambda+\bar p_f-\gamma_f=0\) and the pressure-difference relation
\(\bar p_f-\bar p_g=\gamma_f-\gamma_g\) are useful diagnostics, but they need
not be separate residual equations if the selected primitive variables enforce
\(\bar p_E=-\lambda\) and the capillary closure directly.
The mass-fraction Euler--Lagrange equations provide the composition partition
relation on the constrained tangent space:
\begin{align}
		\rho_\xi
		\left(
		\left.
		\frac{\partial\psi_\xi}{\partial\eta_\xi^\beta}
		\right|_{\mbf F_\xi,\bar\rho_\xi,\eta_\xi^{\kappa\neq\beta},\theta}
		-
		\left.
		\frac{\partial\psi_\xi}{\partial\eta_\xi^N}
		\right|_{\mbf F_\xi,\bar\rho_\xi,\eta_\xi^{\kappa\neq N},\theta}
		\right)
		 & =
		\frac{\pi_\xi^\beta}{\eta_\xi^\beta}
		-\frac{\pi_\xi^N}{\eta_\xi^N},
		\quad \beta=1,\ldots,N-1.
	\label{eq:implementation_pi_difference_system}
\end{align}
These algebraic equations are not optional postprocessing quantities.  They are
part of the nonlinear system that defines \(\phi_\xi\),
\(\eta_\xi^\alpha\), \(\pi_\xi^\alpha\), \(\lambda\), \(\bar p_f\), and
\(\bar p_E\) from the chosen primitive variables.  If all component mass
fractions are solved directly, the unprojected equation with the multiplier
\(\Pi_\xi\) may be used instead, but the projected system is the minimal
independent implementation form.

\subsection{Transfer-potential evolution on the solid frame}
\label{sec:conversion_potential_solid_frame}

Variation of the phase-attached conversion variables first gives a relation
containing the conversion-work potential \(L_\xi^\alpha\).  Subtracting its
reference-component form gives
\begin{equation}
	L_\xi^\alpha-L_\xi^N
	=
	\frac{\pi_\xi^\alpha}
	{\phi_\xi\bar\rho_\xi\eta_\xi^\alpha}
	-
	\frac{\pi_\xi^N}
	{\phi_\xi\bar\rho_\xi\eta_\xi^N}
	=
	\hat\mu_\xi^\alpha-\hat\mu_\xi^N,
	\qquad
	\alpha\ne N .
	\label{eq:implementation_transfer_work_component_difference}
\end{equation}
The neutral potentials are AD material outputs implementing the theory
manuscript's constituent-partial-density derivative
\(\hat\mu_\xi^\alpha
=\partial(\rho_\xi\psi_\xi+\phi_\xi\omega_\xi^+)/
\partial\rho_\xi^\alpha\) with the specified held-fixed fields.  They satisfy
the theory Euler identity
\(\hat\mu_\xi^\alpha
=\psi_\xi+\pi_\xi^\alpha/
(\phi_\xi\bar\rho_\xi\eta_\xi^\alpha)\); neither relation is an additional
finite-element residual.
Choose one skeleton-attached solid phase \(s_\star\) as the transfer reference
and impose
\(\psi_{s_\star}+L_{s_\star}^N=\hat\mu_{s_\star}^N\).
The only evolution equation for the retained transfer potential is then
\begin{align}
	\left.\frac{\partial\tau}{\partial t}\right|_{\mbf X}
	=
	\frac{1}{2}\mbf v_s\cdot\mbf v_s
	+\hat\mu_{s_\star}^N
	-\psi_{s_\star}
	-\frac{\pi_{s_\star}^N}
	{\phi_{s_\star}\bar\rho_{s_\star}\eta_{s_\star}^N}.
	\label{eq:implementation_tau_solid_frame}
\end{align}
Once \(\tau\) is known, the \(P-1\) non-reference phase levels are algebraic.
Using \eqref{eq:fluid_material_derivative_on_solid_frame}, a mobile fluid phase
\(f\ne s_\star\) has
\begin{align}
	L_f^N
	 & =
	\left.\frac{\partial\tau}{\partial t}\right|_{\mbf X}
	+\frac{\mbf W_f}{J\rho_f}\cdot\Grad\tau
	-\frac{1}{2}\mbf v_f\cdot\mbf v_f
	+\frac{\pi_f^N}
	{\phi_f\bar\rho_f\eta_f^N},
	\notag \\
	 & =
	\hat\mu_f^N-\psi_f
	+\left.\frac{\partial\tau}{\partial t}\right|_{\mbf X}
	+\frac{\mbf W_f}{J\rho_f}\cdot\Grad\tau
	-\frac{1}{2}\mbf v_f\cdot\mbf v_f,
	\label{eq:implementation_fluid_transfer_work_recovery}
\end{align}
while another skeleton-attached solid phase \(s\ne s_\star\) has
\begin{align}
	L_s^N
	 & =
	\left.\frac{\partial\tau}{\partial t}\right|_{\mbf X}
	-\frac{1}{2}\mbf v_s\cdot\mbf v_s
	+\frac{\pi_s^N}
	{\phi_s\bar\rho_s\eta_s^N},
	\notag \\
	 & =
	\hat\mu_s^N-\psi_s
	+\left.\frac{\partial\tau}{\partial t}\right|_{\mbf X}
	-\frac{1}{2}\mbf v_s\cdot\mbf v_s.
	\label{eq:implementation_solid_transfer_work_recovery}
\end{align}
These equalities are the direct algebraic closure
\(L_\xi^\alpha=\hat\mu_\xi^\alpha-\psi_\xi
+D_\xi\tau/Dt-\mbf v_\xi\cdot\mbf v_\xi/2\), written on the solid
reference frame.  Thus \(L_\xi^\alpha\) is recovered from the constitutive
potential \(\hat\mu_\xi^\alpha\), the retained field \(\tau\), and the phase
kinematics; it is not an independent constitutive function.

The traditional electrochemical affinity of mechanism \(m\) is the
stoichiometric projection
\begin{equation}
	\mc A_{(m)}
	=
	-\sum_\xi\sum_\alpha
	\mu_\xi^\alpha
	\nu_{\xi (m)}^\alpha .
	\label{eq:implementation_affinity_projection}
\end{equation}
The coefficient supplied by the variational conversion power is related to
this affinity by
\begin{equation}
	-\sum_\xi\sum_\alpha
	\left(
	\psi_\xi+L_\xi^\alpha+z_\xi^\alpha\varphi
	\right)
	\nu_{\xi (m)}^\alpha
	=
	\mc A_{(m)}
	-
	\sum_\xi
	\left(
	\frac{D_\xi\tau}{Dt}
	-\frac{1}{2}\mbf v_\xi\cdot\mbf v_\xi
	\right)
	\sum_\alpha\nu_{\xi (m)}^\alpha .
	\label{eq:implementation_generalized_conversion_affinity_relation}
\end{equation}
The correction vanishes for chemistry that conserves mass separately within
each phase and at full reversible equilibrium.  It must be retained for phase
change, dissolution, precipitation, and pore-filling mechanisms.
Equations~\eqref{eq:implementation_tau_solid_frame}--
\eqref{eq:implementation_generalized_conversion_affinity_relation} are required whenever the
implementation retains \(\tau\) or uses affinity-based kinetics.  For the first
closed problem, \(\tau\) is the single solved transfer field on the solid mesh;
the \(P-1\) levels \(L_\xi^N\) are algebraic variables or AD material
properties, and \eqref{eq:implementation_transfer_work_component_difference}
reconstructs the other component levels.  Their residuals are given in
Section~\ref{sec:conversion_potential_weak_form} so that the affinity and
the generalized conversion coefficient are differentiated consistently by
MOOSE AD.
